% GENERATED FILE. Edit canonical lessons, then run npm run generate:books.
% canonical-source-hash: fnv1a64:8ab8e91e5b4c4481
% canonical-lessons: MR-C62-labels, MR-C62-message, MR-C62-pahila-paath

\chapter{Reading — Labels, a Message, and a First Paragraph}
\label{ch:mr-pahila-paath}
\hlchaptermodality{62}

\begin{chapteropening}
\textbf{By the end of this chapter:} \emph{I can read Marathi off a form, off a short message, and through a whole paragraph, without translating word by word.}

Eighty-three words of Marathi in a row with no new word in them — the first time the track asks for a whole paragraph at once.
\end{chapteropening}

\section[nāv · śahar · bhāṣā · ghar · kholī …]{(reading) — six words, no sentence around them}

\label{lesson:MR-C62-labels}



\begin{quote}
Every word below you have already written. Not one is new.

What is new is that nothing tells you which is coming.
\end{quote}

\begin{tcolorbox}[breakable,skin=enhanced,colback=orange!6,colframe=orange!45!black,boxrule=0.5pt,arc=1mm,left=6pt,right=6pt,top=4pt,bottom=4pt,fonttitle=\bfseries,title={Read this}]
\begin{quote}
\mr{नाव} \mr{शहर} \mr{भाषा} \mr{घर} \mr{खोली} \mr{पाणी}
\end{quote}

Read down the list once. Do not sound out each letter — look at the whole word and let it arrive.

Again, and notice which came at once and which you had to assemble. Both are fine. The difference is what practice removes.
\end{tcolorbox}

\subsection*{You'll want to know: a word with nothing to lean on}
The first three are form labels — the ones you filled in when you wrote your own profile. The last three are things.

A label gives you no sentence to lean on. There is no \textbf{\mr{आहे}} at the end to tell you what is going on, and no verb to wait for. One word, and it carries everything.

That is why six is enough for one sitting.

\subsection*{Your turn}
\begin{itemize}
\raggedright
  \item \emph{Say it:} the six words down the list, once, without stopping
\end{itemize}

\emph{Twice through:}

\begin{tcolorbox}[breakable,colback=teal!4,colframe=teal!35!black,title={Before you move on}]
How many of the six were new? (\textbf{None.})

What was new? (\textbf{Nobody said them first.} You read them.)
\end{tcolorbox}
\section[namaskār mīrā. mājhaṁ nāv aruṇ āhe.]{(reading) — a message with your own sentences in it}

\label{lesson:MR-C62-message}



\begin{quote}
The last lesson gave you single words. This gives each one a sentence to stand in.

You have written every line below. This time somebody wrote it to you.
\end{quote}

\begin{tcolorbox}[breakable,skin=enhanced,colback=orange!6,colframe=orange!45!black,boxrule=0.5pt,arc=1mm,left=6pt,right=6pt,top=4pt,bottom=4pt,fonttitle=\bfseries,title={Read this}]
\begin{quote}
\mr{नमस्कार} \mr{मीरा}. \mr{माझं} \mr{नाव} \mr{अरुण} \mr{आहे}. \mr{भेटून} \mr{आनंद} \mr{झाला}. \mr{मला} \mr{मराठी} \mr{येते}. \mr{मी} \mr{मराठी} \mr{वाचतो}. \mr{धन्यवाद}.
\end{quote}

Read it through once, without stopping.

Again — and notice you did not assemble it line by line the second time. You took it in.
\end{tcolorbox}

\subsection*{You'll want to know: the same sentences, the other way round}
These are the lines from your own six-field profile and your own message, in the order you wrote them.

Writing them, you chose each one and had all the time you wanted. Reading them, they arrive in somebody else's order and at their pace. That is a different skill, and it is why this is a separate lesson rather than a review.

Six lines is not much. Meeting them from the other side is the whole exercise.

\subsection*{Your turn}
\begin{itemize}
\raggedright
  \item \emph{Say it:} the six lines, once, in order, without stopping
\end{itemize}

\emph{Twice through:}

\begin{tcolorbox}[breakable,colback=teal!4,colframe=teal!35!black,title={Before you move on}]
Who is the message addressed to? (\textbf{\mr{मीरा}.})

How many of the six lines were new? (\textbf{None.})
\end{tcolorbox}
\section[mājhaṁ nāv aruṇ āhe. malā marāṭhī yete.]{(reading) — the first paragraph}

\label{lesson:MR-C62-pahila-paath}



\begin{quote}
Six words, then six lines. Now a whole paragraph.

Eighty-three words, and the rule has not changed: not one of them is new.
\end{quote}

\begin{tcolorbox}[breakable,skin=enhanced,colback=orange!6,colframe=orange!45!black,boxrule=0.5pt,arc=1mm,left=6pt,right=6pt,top=4pt,bottom=4pt,fonttitle=\bfseries,title={Read this}]
\begin{quote}
\mr{माझं} \mr{नाव} \mr{अरुण} \mr{आहे}. \mr{मला} \mr{मराठी} \mr{येते}. \mr{माझं} \mr{घर} \mr{मोठं} \mr{आहे}. \mr{घरामध्ये} \mr{एक} \mr{खोली} \mr{आहे}. \mr{खोली} \mr{लहान} \mr{आहे} \mr{पण} \mr{चांगली} \mr{आहे}. \mr{मी} \mr{चहा} \mr{पितो}. \mr{मला} \mr{दूध} \mr{सुद्धा} \mr{आवडतं}. \mr{मला} \mr{भाकरी} \mr{आवडते}. \mr{माझा} \mr{भाऊ} \mr{उंच} \mr{आहे}. \mr{माझा} \mr{मित्र} \mr{जवळ} \mr{आहे}. \mr{मित्राचे} \mr{नाव} \mr{मीरा} \mr{आहे}. \mr{मीरा} \mr{हुशार} \mr{आहे} \mr{आणि} \mr{खूप} \mr{वाचते}. \mr{मीरा} \mr{मला} \mr{मदत} \mr{करते}. \mr{मीरा} \mr{माझ्याबरोबर} \mr{चालते}. \mr{मी} \mr{मराठी} \mr{वाचतो} \mr{आणि} \mr{लिहितो}. \mr{मीरा} \mr{म्हणते} \mr{की} \mr{मराठी} \mr{सोपं} \mr{आहे}. \mr{मला} \mr{पण} \mr{सोपं} \mr{वाटतं}, \mr{कारण} \mr{मी} \mr{बोलतो}. \mr{मीरा} \mr{विचारते}, \mr{किती}? \mr{मी} \mr{म्हणतो}, \mr{दहा}. \mr{उद्या} \mr{भेटू}. \mr{धन्यवाद}.
\end{quote}

Read it once without stopping. Do not translate as you go — let the sentences arrive.

Now read it again, and notice how little work it took.
\end{tcolorbox}

\subsection*{You'll want to know: what is holding it together}
\textbf{\mr{आणि}} adds. \textbf{\mr{पण}} holds two true things apart. \textbf{\mr{कारण}} gives a reason. \textbf{\mr{की}} hangs a whole sentence off another one. \textbf{\mr{सुद्धा}} adds a thing to a list already made. Each of those was a chapter, and each was small.

Notice also that Marathi waits. \textbf{\mr{मीरा} \mr{म्हणते} \mr{की} \mr{मराठी} \mr{सोपं} \mr{आहे}} — you do not find out what Mira says until the end, and you held the beginning without being asked to.

This is the first whole paragraph the track has asked you to hold at once, and it works because every piece was paid for.

\subsection*{Your turn}
\begin{itemize}
\raggedright
  \item \emph{Say it:} the paragraph aloud, once, without stopping
\end{itemize}

\emph{Twice through:}

\begin{tcolorbox}[breakable,colback=teal!4,colframe=teal!35!black,title={Before you move on}]
How many words in that paragraph were new to you? (\textbf{None.})

And what made it readable? (Everything in it had already been \textbf{paid for}, one chapter at a time.)
\end{tcolorbox}
